% Paper III -- Corrected Table III (replaces tab:udg in Paper I)
\begin{table}[h]
\caption{DEV gravitational slip predictions for known UDGs,
recomputed with published photometric parameters and using the
argument $x \equiv g_N/\aO$ as in Paper~I, Sec.~III.
Two conventions are tabulated: $\eta-1$ evaluated at the
Newtonian $g_N$ (column~5, the original Paper~I convention) and
at the observed $g_{\rm obs}=\nu(g_N/\aO)\,g_N$ (column~6,
the convention adopted in Paper~III).}
\label{tab:udg}
\begin{ruledtabular}
\begin{tabular}{lcccccc}
System & $M_\star$ & $r_{\rm eff}$ & $g_N/\aO$ & $g_{\rm obs}/\aO$ & $(\eta-1)_{g_N}$\,[\%] & $(\eta-1)_{g_{\rm obs}}$\,[\%] \\
       & ($10^8 M_\odot$) & (kpc) & & & & \\
\hline
NGC1052-DF2 & 2.0 & 2.2 & 0.0480 & 0.2218 & 2.23 & 0.96 \\
NGC1052-DF4 & 1.5 & 2.0 & 0.0436 & 0.2110 & 2.34 & 0.99 \\
DF44 & 3.0 & 4.3 & 0.0189 & 0.1379 & 3.61 & 1.26 \\
DGSAT-I & 3.0 & 4.7 & 0.0158 & 0.1261 & 3.95 & 1.33 \\
VCC1287 & 1.0 & 2.8 & 0.0148 & 0.1222 & 4.08 & 1.35 \\
DF17 & 2.0 & 3.5 & 0.0190 & 0.1384 & 3.60 & 1.26 \\
\end{tabular}
\end{ruledtabular}
{\small Note: the previous version of this table (dev\_paper\_I\_corrected.tex)
used $r_{\rm eff}=11.8$\,kpc for DGSAT-I, which yielded
$g_N/\aO=0.0025$ and $\eta-1=9.9\%$. With the photometric value
$r_{\rm eff}=4.7$\,kpc (Janowiecki+15), $g_N/\aO=0.016$ and the
prediction collapses to $\eta-1\simeq 4.0\%$ (or $1.3\%$ in the
$g_{\rm obs}$ convention). All six UDGs lie in the deep-MOND regime
and predict $\eta-1\sim 1$--$5\%$, near the Euclid threshold.}
\end{table}